\section{From infinitely many midpoints to diamond distortion}\label{met:section}
A symmetric midpoint estimate becomes a graph obstruction because each edge
of a countably branching diamond has infinitely many replacement midpoints.
Finite-dimensional compactness lets us compare two of them while making
any prescribed finite set of coordinates negligible. We give the transfer
with all scaling factors and deduce the principal squared-loss bound.
Section~\ref{met:transfer-section} develops further ways to accumulate
midpoint estimates. These are instances of the midpoint method for countably
branching diamonds; see \cite[Section~4]{BaudierEtAl2017}. The quantitative
constants below come from the estimates stated here.

For $x\in X$ and $R\ge0$, write
$\mathcal L_R(x)=\{z\in X:\|x+z\|\le R,\ \|x-z\|\le R\}$
for the symmetric $R$-lens at $x$. Its members are deviations from the
common midpoint, rather than the two endpoints themselves.

\subsection{The graph and its slopes}\label{met:graph-definition}
Let $\mathcal D_0$ consist of two vertices joined by an edge of length one.
To obtain $\mathcal D_{j+1}$, replace every edge $[a,b]$ of $\mathcal D_j$
by paths $[a,m_i,b]$, indexed by $i\in\mathbb N$, whose middle vertices
are distinct and new. Both new edges have half the length of their parent.
Let $d_k$ be the shortest-path metric on the vertices of $\mathcal D_k$.
We call an edge present at stage $j$ a \emph{stage-$j$ edge}; it need no
longer be an edge at stage $k>j$.

Distances between old vertices are unchanged at a replacement. Indeed,
an old path can be replaced by paths of the same total length. Conversely,
a path between old vertices consists of two-edge portions through new
vertices; each portion either returns to the same old vertex or traverses
the endpoints of one old edge. Compressing these portions gives an old
walk of no greater length. It follows that a stage-$j$ edge has endpoint
distance $\ell_j=2^{-j}$ in $\mathcal D_k$. Its distinct replacement
midpoints have mutual distance $\ell_j$, and distance $\ell_j/2$ from
either endpoint. At their introduction the assertion is immediate because
two distinct middle vertices are nonadjacent, and every incident edge has
length $\ell_j/2$; subsequent replacements preserve it.

\begin{figure}[!bp]
\centering
\begin{tikzpicture}[scale=0.85,every node/.style={font=\small}]
 \draw[thick] (0,0)--node[above] {$\ell$}(2,0);
 \fill (0,0) circle (2pt) node[left] {$u$};
 \fill (2,0) circle (2pt) node[right] {$v$};
 \draw[-{Stealth},thick] (2.8,0)--(4,0);
 \coordinate (u) at (5,0); \coordinate (v) at (9,0);
 \foreach \h/\lab in {1.4/1,0.45/2,-1.4/j}{
   \draw[thick] (u)--(7,\h)--(v);
   \fill (7,\h) circle (2pt) node[above] {$w_{\lab}$};
 }
 \node at (5.75,1.0) {$\ell/2$};
 \node at (8.25,1.0) {$\ell/2$};
 \node at (7,-0.5) {$\vdots$};
 \fill (u) circle (2pt) node[left] {$u$};
 \fill (v) circle (2pt) node[right] {$v$};
 \node at (1,-1.9) {one old edge};
 \node at (7,-1.9) {countably many two-edge paths};
\end{tikzpicture}
\caption{One normalized diamond replacement, with only three of the
countably many paths shown. Each new edge has length $\ell/2$.
The old endpoints remain at distance $\ell$, and any two distinct new
middle vertices also have distance $\ell$. The drawing is schematic.}
\label{fig:diamond}
\end{figure}


The normalized metric $d_k$ defined above differs by the global factor
$2^{-k}$ from the convention
in which every edge of the final graph has length one. In that convention
old-vertex distances double at each replacement and a stage-$j$ edge has
represented length $2^{k-j}$. Both conventions give exactly the same
embedding distortion.

Suppose $f:\mathcal D_k\to X$, $s>0$ and $D\ge1$ satisfy
\begin{equation}\label{met:all-pairs}
 s d_k(u,v)\le\|f(u)-f(v)\|\le Ds d_k(u,v)
 \qquad(u,v\in\mathcal D_k).
\end{equation}
The lower inequality is required for all pairs, including distinct
replacement midpoints. Dividing $f$ by $s$ gives the noncontracting
normalization $d_k\le\|f(u)-f(v)\|\le Dd_k$. Dividing instead by $Ds$
gives lower factor $\alpha=1/D$ and upper factor one. We use both
normalizations and state which one is in force.

Orient the original edge and orient every replacement path consistently.
For a parent edge $[a,b]$ of represented length $\ell$, define
\begin{equation}\label{met:slopes}
 x=\frac{f(b)-f(a)}{\ell},\qquad
 z_i=\frac{2}{\ell}\left(f(m_i)-\frac{f(a)+f(b)}2\right).
\end{equation}
Its child slopes are $x+z_i$ and $x-z_i$. Under lower factor $\alpha$,
\begin{equation}\label{met:sibling-separation}
 \|z_i-z_j\|\ge2\alpha\qquad(i\ne j).
\end{equation}
Every slope lies between $\alpha$ and the upper embedding factor. The
sequence $(z_i)$ is bounded, since each $z_i$ is half the difference of
two child slopes.

\subsection{The finite-dimensional pair selection}\label{met:pairing-section}
The next observation records the approximation errors that will be used
throughout. It requires bounded finite-rank projections, rather than
contractivity of the complementary projections.

\begin{lemma}\label{met:finite-pair}
Let $X$ be a Banach space, $R>0$, and let $(z_i)_{i\in\mathbb N}$ satisfy
$\|x\pm z_i\|\le R$ with $\|z_i-z_j\|\ge\varepsilon>0$ for $i\ne j$.
Every half-difference $w=(z_i-z_j)/2$ with $i\ne j$ satisfies
$\|w\|\ge\varepsilon/2$ and $\|x\pm w\|\le R$.
Fix $\eta>0$. The following two selections are available.
\begin{enumerate}
\item Suppose $p$ has finite support, $\|p-x\|<\eta$, and $P$ is a bounded
finite-rank coordinate projection satisfying $Pp=p$. There are distinct
$i,j$ for which, with $w=(z_i-z_j)/2$ and $y=w-Pw$,
\begin{equation}\label{met:pair-errors}
 \|Pw\|<\eta,\quad
 \|y\|\ge\varepsilon/2-\eta,\quad
 \|p\|\ge\|x\|-\eta,\quad
\|p\pm y\|\le R+2\eta.
\end{equation}
\item If $F$ is a closed finite-codimensional subspace, there are distinct
$i,j$ whose half-difference satisfies
$\operatorname{dist}(w,F)<\eta$.
This selection requires no coordinate structure on $X$.
\end{enumerate}
\end{lemma}
\begin{proof}
For the common endpoint bound use
\[
 x+w=\tfrac12((x+z_i)+(x-z_j)),\qquad
x-w=\tfrac12((x-z_i)+(x+z_j)).
\]
In the first selection, the bounded sequence $(Pz_i)$ has two distinct
members at distance less than $2\eta$, because bounded subsets of a
finite-dimensional space are totally bounded. This gives the first
inequality in~\eqref{met:pair-errors} and hence the second.
Replacing $x$ by $p$ and $w$ by $w-Pw$ costs less than $2\eta$.
The last assertion follows by applying the same argument to the bounded
images of $z_i$ in the finite-dimensional quotient $X/F$.
\end{proof}

\subsection{The quadratic loss from the pooled atomic estimate}\label{met:joint-section}
For the segment forest space, the pooled atomic estimate has the form
\begin{equation}\label{met:joint-input}
 \|Q_Hy\|\le2\sqrt{R^2-\|p\|^2},
 \qquad P_Hp=p,\quad \|p\pm y\|\le R.
\end{equation}
Here $H$ is finite and ancestral, $Q_H=I-P_H$, and the displacement
$y$ is arbitrary. The estimate is supplied by the joint-budget
theorem in the companion on midpoint lenses
\cite[Theorem~1.1]{SegmentLenses}. Its target is the
projected tail $Q_Hy$; an individual endpoint-law estimate is not being
used with an arbitrary head.

The same companion establishes, for both the finite-forest dual and
the infinite-height closed coordinate span, density of finite coordinate
vectors and contractivity of $P_H$ and $Q_H$ for ancestral $H$.
Thus $P_Hx$ can approximate any center $x$: take a finite coordinate
approximation and let $H$ contain its support and all its ancestors.
These approximation properties refer to each space with its own norm.

\begin{lemma}\label{met:joint-lens-lemma}
In the finite-height segment forest dual, every infinite
$\varepsilon$-separated symmetric $R$-lens satisfies
\begin{equation}\label{met:joint-lens}
 \|x\|^2+\varepsilon^2/16\le R^2.
\end{equation}
The same lens conclusion holds for the infinite-height closed coordinate
span with its own segment-dual norm and finite-head approximation theorem.
\end{lemma}
\begin{proof}
Fix $0<\eta<\varepsilon/2$ and choose $H$ so that
$p=P_Hx$ has $\|p-x\|<\eta$. Choose two deviations with
$w=(z_i-z_j)/2$ satisfying $\|P_Hw\|<\eta$.
Then $\|Q_Hw\|\ge\varepsilon/2-\eta$ and
$\|p\pm w\|\le R+\eta$, because $\|x\pm w\|\le R$.
Apply~\eqref{met:joint-input} to the arbitrary displacement $w$:
\[
 (\varepsilon/2-\eta)^2
 \le4\bigl((R+\eta)^2-\|p\|^2\bigr).
\]
Letting $\eta\downarrow0$ proves~\eqref{met:joint-lens}. The same
argument applies with the infinite-height coordinate span's own version
of~\eqref{met:joint-input}.
\end{proof}

\begin{proof}[Proof of Theorem~\ref{met:intro-main}]
Normalize to lower factor one and upper factor $D$.
Let $R_j$ be the supremum of stage-$j$ slope norms. Every parent
slope is the center of a $2$-separated lens whose endpoints have norms
at most $R_{j+1}$. Thus $\|x_e\|^2+1/4\le R_{j+1}^2$ for every
parent. Only now take the supremum: $R_j^2+1/4\le R_{j+1}^2$.
Since $R_0\ge1$ and $R_k\le D$, summing gives $D^2\ge1+k/4$.
No supremum need be attained. The infinite-height statement uses exactly
the same scalar argument with that space's proved version of
\eqref{met:joint-input}; no identification with the full infinite-height
dual is required.
\end{proof}

The $1/4$ coefficient comes from retaining the joint quadratic budget
in~\eqref{met:joint-input}. The direct graph-energy arguments follow next. General transfer rules
and alternative forest estimates appear in Section~\ref{met:transfer-section}
and Appendix~\ref{met:forest-variants}.
